\section{Introduction}
Travel demand during extended holidays can differ substantially from an ordinary workday. Taiwan's Freeway Bureau reported a daily average of 124.0 million vehicle-kilometers across the national freeway network from the first through sixth days of the 2026 Lunar New Year, approximately 1.3 times the annual workday average. Over the same six days, approximately 77,000 vehicles per day passed through the Hsuehshan Tunnel on National Freeway No.~5 \cite{fwb2026holiday}. The Bureau also publishes traffic controls that can vary by interchange, direction, and time, including ramp closures and high-occupancy vehicle restrictions \cite{fwb2026controls}. A traveler deciding whether to leave now or wait needs both short-term traffic information and the controls applicable to the trip.

We study short-horizon speed forecasting on ten directional segments of National Freeway No.~5. The proposed agent combines predictions from electronic toll collection (ETC) data with retrieved notices. Its output is designed to show a departure option, the supporting predicted speeds and risks, the applicable controls, and the limits of the forecast. Notice retrieval provides context for the advice. Establishing whether a control causes a change in congestion would require a separate analysis.

The study contributes a route-specific evaluation of speed and congestion-onset forecasts, including a separate view of holiday performance. It also specifies a tool-based agent architecture that checks generated numeric claims against the forecast output, and a data-cleaning rule that distinguishes a zero-speed record with no observed vehicles from a valid observation. The preliminary experiment evaluates the forecasting component. The quality and benefit of the departure advice are future evaluation targets.
